\medskip

Dans l'espace rapporté à un repère orthonormé \Oijk, on considère:

\begin{itemize}
\item[$\bullet~~$] la droite $\mathcal{D}$ passant par le point A(2~;~4~;~0) et dont un vecteur directeur est $\vect{u}\begin{pmatrix}1\\2\\0\end{pmatrix}$ ;
\item[$\bullet~~$] la droite $\mathcal{D}'$ dont une représentation paramétrique est : $\left\{\begin{array}{l !{=} l}
x&3\\
y& 3 + t\\
z&3 + t
\end{array}\right. ,\: t \in \R$.
\end{itemize}



\begin{enumerate}
\item 

	\begin{enumerate}
		\item 
La droite $\mathcal{D}'$ a pour vecteur directeur $\vect{u'}$ de  coordonnées 
$\begin{pmatrix}
0\\1\\1
\end{pmatrix}$

		\item Les vecteurs $\vect{u}\begin{pmatrix}1\\2\\0\end{pmatrix}$ et  $\vect{u'}\begin{pmatrix}0\\1\\1\end{pmatrix}$ ne sont pas colinéaires, donc les droites $\mathcal{D}$ et $\mathcal{D}'$ ne sont pas parallèles.

		\item La droite $\mathcal{D}$ est l'ensemble des points M$\left (x\,; y\;;z\right )$ tels que $\vectt{AM}$ et $\vect{u}$ soient colinéaires, c'est-à-dire $\vectt{AM}= t.\vect{u}$ avec $t\in\R$, soit:	
$\left\{\begin{array}{l !{=} l}
x-2&t\\
y-4& 2 t\\
z-0&0t
\end{array}\right. ,\: t \in \R$
		
		
La droite $\mathcal{D}$ a pour représentation paramétrique
$\left\{\begin{array}{l !{=} l}
x&2+t\\
y&4+ 2 t\\
z&0
\end{array}\right. ,\: t \in \R$

	\end{enumerate}

\end{enumerate}

On admet dans la suite de cet exercice qu'il existe une unique droite $\Delta$ perpendiculaire aux droites $\mathcal{D}$ et $\mathcal{D}'$. Cette droite $\Delta$ coupe chacune des droites $\mathcal{D}$ et $\mathcal{D}'$. On appellera M le point d'intersection de $\Delta$ et $\mathcal{D}$, et M$'$ le point d'intersection de $\Delta$ et $\mathcal{D}'$.

On se propose de déterminer la distance MM$'$ appelée \og distance entre les droites $\mathcal{D}$ et $\mathcal{D}'$ \fg.

\begin{enumerate}[resume]
\item Soit le vecteur $\vect{v}\begin{pmatrix}2\\- 1\\1\end{pmatrix}$.

\begin{list}{\textbullet}{}
\item $\vect{v}\cdot\vect{u}=2\times 1 + (-1)\times 2 +1\times 0 = 0$ donc $\vect{v}\perp \vect{u}$.
\item $\vect{v}\cdot\vect{u'}=2\times 0 + (-1)\times 1 +1\times 1 = 0$ donc $\vect{v}\perp \vect{u'}$.
\end{list}

Le vecteur $\vect{v}$ est donc orthogonal aux deux vecteurs directeurs de $\mathcal{D}$ et $\mathcal{D}'$, donc c'est un vecteur directeur de leur perpendiculaire commune $\Delta$.

\item On note $\mathcal{P}$ le plan contenant les droites $\mathcal{D}$ et $\Delta$, c'est-à-dire le plan passant par le point A et de vecteurs directeurs $\vect{u}$ et $\vect{v}$.
	\begin{enumerate}
		\item Soit le vecteur $\vect{n}\begin{pmatrix}2\\-1\\-5 \end{pmatrix}$.
		
\begin{list}{\textbullet}{}
\item $\vect{n}\cdot\vect{u}=2\times 1 + (-1)\times 2 +(-5)\times 0 = 0$ donc $\vect{n}\perp \vect{u}$.
\item $\vect{n}\cdot\vect{v}=2\times 2 + (-1)\times (-1) +(-5)\times 1 = 0$ donc $\vect{n}\perp \vect{v}$.
\end{list}		

Le vecteur $\vect{n}$ est donc orthogonal à deux vecteurs directeurs du plan $\mathcal{P}$, donc le vecteur $\vect{n}$ est un vecteur normal au plan $\mathcal{P}$.
		
		\item 
Le plan $\mathcal{P}$ est le plan ayant $\vect n$ comme vecteur normal et passant par A, donc c'est  l'ensemble des points P$\left (x\,; y\;;z\right )$ tels que 
$\vectt{AP} \perp \vect n$.

$\vectt{AP}$ a pour coordonnées 
$\begin{pmatrix}
x-2\\y-4\\z
\end{pmatrix}$

$\begin{gathered}\vectt{AP} \perp \vect n
\iff \vectt{AP} \cdot \vect n
\iff (x-2)\times 2 + (y-4)\times (-1) + z\times (-5)=0\\
\phantom{\vectt{AP} \perp \vect n \iff \vectt{AP} \cdot \vect n}
\iff 2x-4-y+4 -5z=0
\iff 2x-y-5z=0\end{gathered}$

Donc une équation du plan $\mathcal{P}$ est : $2x - y - 5z = 0$.		
		
		\item On rappelle que M$'$ est le point d'intersection des droites $\Delta$ et $\mathcal{D}'$.
		


Le plan $\mathcal{P}$ contient la droite $\Delta$ et M$'$ est un point de la droite $\Delta$; donc M$'$ est un point du plan $\mathcal{P}$.
M$'$ appartient à $\mathcal{D}'$ et à $\mathcal{P}$ donc M$'$ est  le point d'intersection de la droite $\mathcal{D}'$ et du plan $\mathcal{P}$.
		


Les coordonnées $\left (x\,; y\;;z\right )$ de M$'$ vérifient donc le système
$\left \{
\begin{array}{r !{=} l}
x & 3\\
y & 3+t\\
z & 3+t\\
2x-y-5z & 0
\end{array}
\right .$

$2x-y-5z=0$ devient $2\times 3 -(3+t) -5(3+t)=0$ soit $6-3-t-15-5t=0$ ou encore $-12=6t$ ce qui donne $t=-2$.

$x=3$; $y=3+t=3-2=1$ et $z=3+t=1$

Les coordonnées du point M$'$ sont donc (3~;~1~;~1).

	\end{enumerate}
\item 
	\begin{enumerate}
		\item La droite $\Delta$ a pour vecteur directeur $\vect v$ et passe par le point M$'$,  donc elle a pour représentation paramétrique:
$\left \{
\begin{array}{r !{=} l}
x & 3+2t'\\
y & 1-t'\\
z & 1+t'
\end{array}
\right . \quad t'\in\R$.		
		
		\item 
Le point M est le point d'intersection de $\mathcal{D}$ et de $\Delta$ donc ses coordonnées $\left (x\,; y\;;z\right )$ sont solutions du système:
$\left \{
\begin{array}{r !{=} l}
x&2+t\\
y&4+ 2 t\\
z&0\\
x & 3+2t'\\
y & 1-t'\\
z & 1+t'
\end{array}
\right . $
		
On a donc 	
$\left \{
\begin{array}{r !{=} l}
2+t & 3+2t'\\
4+2t & 1-t'\\
0 & 1+t'
\end{array}
\right . $
ou encore
$\left \{
\begin{array}{r !{=} l}
t & -1\\
-1 & t'
\end{array}
\right . $

$\left \{
\begin{array}{r !{=} l}
x&2+t\\
y&4+ 2 t\\
z&0
\end{array}
\right . $	
donc
$\left \{
\begin{array}{r !{=} l}
x&1\\
y&2\\
z&0
\end{array}
\right . $

Donc  le point M a pour coordonnées (1~;~2~;~0).

		\item 
$\begin{gathered}\text{MM}' = \ds\sqrt{(x_{\text{M}'} - x_{\text{M}})^2+(y_{\text{M}'} - y_{\text{M}})^2+ (z_{\text{M}'} - z_{\text{M}})^2}\\
\phantom{\text{MM}'}
=\ds\sqrt{(3-1)^2+(1-2)^2+(1-0)^2}
=\ds\sqrt{4+1+1}
=\ds\sqrt{6}\end{gathered}$
	\end{enumerate}
	
\item On considère la droite $d$ de représentation paramétrique 
$\left\{\begin{array}{l !{=} r}
x & 5t\\y & 2 + 5t \\z &1 + t\end{array}\right.$ avec $t \in \R$.
	\begin{enumerate}
		\item 
On cherche	l'intersection de la droite $d$ et du plan $\mathcal{P}$, et pour cela on résout le système:
$\left\{\begin{array}{r !{=} r}
x & 5t\\y & 2 + 5t \\z &1 + t \\ 2x-y-5z & 0\end{array}\right.$

On a donc $2\times (5t) -(2+5t) -5(1+t)=0$ soit $10t -2-5t-5-5t=0$ ou  $0t=7$.

Le système n'a pas de solution donc la droite $d$ est strictement parallèle à $\mathcal{P}$.
		
		\item On note $\ell$ la distance d'un point N de la droite $d$ au plan $\mathcal{P}$ . 
		
On veut exprimer le volume $V$ du  tétraèdre ANMM$'$ en fonction de $\ell$.

$V = \dfrac13 \times B \times h$ où $B$ désigne l'aire d'une base et $h$ la hauteur relative à cette base.

Les points A, M et M$'$ appartiennent au plan $\mathcal{P}$ donc le triangle AMM$'$ forme une base du tétraèdre dont la hauteur est la distance du point N au plan $\mathcal{P}$, c'est-à-dire $\ell$.

La droite $\Delta$ est la perpendiculaire commune à $\mathcal{D}$ et $\mathcal{D}'$ donc $\mathcal{D}$ est perpendiculaire à $\Delta$.

A appartient à $\mathcal{D}$, M est le point d'intersection de $\mathcal{D}$ et de $\Delta$, M$'$ appartient à $\Delta$, et les droites $\mathcal{D}$ et $\Delta$ sont perpendiculaires. On peut en déduire que le triangle AMM$'$ est rectangle en M.

$\begin{gathered}\text{AM} = \ds\sqrt{(x_{\text{M}} - x_{\text{A}})^2+(y_{\text{M}} - y_{\text{A}})^2+ (z_{\text{M}} - z_{\text{A}})^2}\\
\phantom{\text{AM}}
= \ds\sqrt{(1-2)^2 + (2-4)^2 + (0-0)^2} = \ds\sqrt{5}\end{gathered}$

L'aire de la base, c'est-à-dire AMM$'$ vaut $\dfrac{1}{2}\times\text{AM}\times \text{MM}' = \dfrac{\sqrt{5}\times \sqrt{6}}{2}=\dfrac{\sqrt{30}}{2}$.

Le volume du tétraèdre vaut donc $\dfrac{1}{3}\times \dfrac{\sqrt{30}}{2}\times \ell = \dfrac{\ell\sqrt{30}}{6}$.

		\item N$_1$ et N$_2$ sont deux points quelconques de la droite $d$.
		
La droite $d$ est strictement parallèle au plan $\mathcal{P}$ donc les distances de N$_1$ et N$_2$ au plan $\mathcal{P}$ sont égales.

Les bases des  tétraèdres AN$_1$MM$'$ et AN$_2$MM$'$ sont identiques (le triangle AMM$'$).

Donc les tétraèdres AN$_1$MM$'$ et AN$_2$MM$'$ ont le même volume.
		
	\end{enumerate}
\end{enumerate}
